\chapter{Who Decides What Language Is Correct?}
\section{An Exercise Book Marked in Red}
Pick up something you may know well: a Chinese-language exercise book just handed back to you.

Open it to last week's composition. You wrote about going to the market with your grandmother at the weekend. One sentence reads, in regional phrasing, ``Jinzhao the vegetables were quite cheap, and Grandmother was jiaoguan happy.'' The teacher has circled jinzhao and jiaoguan in red and written, ``Dialect: nonstandard.'' You lost two marks.

You feel this is unfair. Jinzhao means ``today'', and jiaoguan means ``very''. Grandmother says them daily, and everyone at the market understands. How did they become errors in an exercise book? More puzzling still, writing jintian for ``today'' is correct, while writing jinzhao for the same thing is wrong. Who decided these two fates?

Turn a few pages and things become subtler. The teacher has not circled every ``odd word''. You wrote geili, ``awesome'', without a red circle, and dianzan, ``give a like'', without one either. A classmate even wrote ``feeling emo'' and received only a wavy underline. These words did not exist a few years ago, and most dictionaries do not provide standard definitions. Why do they pass when jinzhao does not?

This exercise book contains the chapter's entire problem. Language is not a fixed instruction manual; it changes every day. Yet schools, examination papers, and broadcasts do draw a line called ``correctness'', awarding marks on one side and deducting them on the other. Who draws that line, and on what grounds? What do people outside it---dialect speakers, people with accents, users of new words---lose, and what does that loss mean?

There is a deeper layer. The red circle appears to correct your vocabulary, but it also quietly edits your relationship with Grandmother. It suggests her speech is inferior, unpresentable, and embarrassing in an examination. The circle is only millimetres wide, but it encloses a person's origins.

That is our question: who has the authority to decide what language is ``correct''? Where does this power come from, and who uses it for what?

\section{The Weight of a Wooden Sign}
Come into another scene. Historical records show that in eighteenth- and nineteenth-century Wales, the mountainous land west of England, many village schools used a particular punishment. From those records, let us reconstruct one school day.

The place is a one-room village school in northern Wales. Stone walls, wooden floor, misted windows. More than thirty children crowd onto long benches; wet wool and coal smoke scent the air. At the front, the teacher leads an English reading lesson. The children repeat, one after another.

A little girl called Mair forgets herself at break. Her neighbour's pencil rolls onto the floor. Bending to pick it up, she blurts something in Welsh, the language she speaks every day at home, in church, and at market. The classroom falls silent.

The teacher approaches carrying a wooden sign, about palm-sized, hanging from a cord. Two words are carved on it: ``Welsh Not'', meaning no speaking Welsh. The rule is this: anyone caught speaking Welsh must wear the sign. Its wearer can ``catch'' another Welsh-speaking pupil and pass it on. Whoever wears it at the end of the school day is beaten.

For the rest of the day, a silent war takes place. Children listen for one another's slips. Mair keeps her mouth tightly shut, afraid to speak even when answering a question. Her English is not yet good enough, and she fears Welsh will escape in a moment of urgency. That afternoon, a boy speaks his mother tongue while helping his younger sister pick something up. Relieved, Mair passes him the sign. When the bell rings, it hangs around a still younger child's neck. The child cries, but rules are rules.

Stay in this classroom a little longer and notice several details.

First, the punishment is enforced not only by the teacher but by the children themselves. The sign's cleverest feature makes the governed monitor one another. Once a linguistic rule exists, maintaining it becomes everyone's labour, including its victims'.

Second, these children still speak Welsh at home. Most of their parents know only Welsh. School language and the language of generations at home have been artificially cut apart: one half is ``correct''; the other gets you beaten.

Third, this is not one teacher's eccentricity. In the middle of the nineteenth century, the British government investigated education in Wales. Its 1847 report described Welsh as an obstacle to education and a breeding ground for backwardness and moral problems. Speaking English became a passport to advancement; speaking Welsh became a sign of having no prospects. The wooden sign was that larger judgement taking shape around a child's neck.

Similar stories occurred far beyond Wales. In France during the same period, public-school children speaking Breton, Basque, or southern dialects could likewise be made to wear signs, stand as punishment, or copy lines. In Chinese classrooms, generation after generation learned ``Speak Putonghua; be a civilised person''. In Japanese schools, children from Okinawa or the rural northeast struggled to eliminate their accents to avoid mockery. Places and materials differed, but the wooden sign remained.

Remember this sign. Only palm-sized, it carries everything we must weigh in this chapter: when one language is declared ``correct'', what happens to people who live in another?

\section{Whose Correctness Is It?}
Lay out the conflict before giving a verdict.

Return to your exercise book. The teacher's reasons for circling jinzhao are quite complete. Exams require standard written language; regional vocabulary may confuse markers from elsewhere. Schools must teach a common way of writing. Otherwise students nationwide would each write differently, making examination marking impossible. This is not merely the teacher's personal prejudice but a premise underlying schooling, examinations, and publishing.

Your objection is reasonable too. Language lives. Geili and dianzan were once ``nonstandard'' novelties and now appear in mainstream media. Jinzhao has been used in Wu Chinese for over a thousand years: why call it an error? If ``correct'' means ``everyone understands'', everyone at the market understands jinzhao. If it means ``in the dictionary'', dictionaries follow language, not the other way around.

Notice that the real question is not ``Are dialects good?'' It is:

\textbf{Where does the judgement ``correct'' come from?}

Does language carry right and wrong within itself, like a mathematical answer everyone calculates identically? Or does one group have authority to announce right and wrong while everyone else obeys? If so, who are they, and what gives them that authority? Are they describing language as it exists, or using language for something else, such as sorting who deserves good jobs, good schools, and good prospects?

Look one layer deeper. ``Standard language'' sounds neutral, like a shared ruler. But rulers do not fall from the sky. China's official definition of Putonghua specifies Beijing pronunciation as the pronunciation standard, northern speech as its dialectal basis, and exemplary modern vernacular writing as its grammatical norm. Study that definition for three seconds. It does not say ``the best pronunciation''; it chooses the speech of \textbf{a particular place}. Why Beijing rather than Suzhou? Why does the standard always happen to reside in powerful cities, cultivated classes, and the mouths of those who write examinations? Has any standard ever grown from markets and villages?

Before reading further, answer for yourself: when the school circles jinzhao, is it maintaining a rule or exercising power? Or are those the same thing?

\section{Two Voices: Standards and Living Speech}
\textbf{Voice one: without a standard, everyone suffers.}

Let me present this case fully. It deserves care because advocates of standard language are easily portrayed as villains strangling diversity, which is unfair.

The first argument is efficiency. In a country with thousands of dialects, producing military call-up orders, medicine instructions, railway rules, and court judgements in every variety would cost astronomical sums. A common language allows strangers to cooperate cheaply, providing a foundation for the modern state. The second argument is fairness, often overlooked. Revisit the Welsh classroom from another angle. Many nineteenth-century Welsh mining parents \textbf{actively} wanted their children to learn English, because jobs outside the mines, urban positions, and government paperwork all used it. Not knowing the standard language did not mean ``preserving yourself'' but being locked in place. For disadvantaged children, standard language was often their only ticket to tangible improvement, not an instrument of oppression. The third argument is historical comparison. Before writing was standardised, a word could have ten spellings, contracts could be evaded through verbal quibbles, and official orders changed meaning on their way to localities. Only after spelling and grammar were standardised did literacy education, dictionaries, and printed publishing truly become possible.

There is hard evidence behind this reasoning. In 1928, Turkey replaced the Arabic alphabet it had used for centuries with a Latin alphabet and then launched mass literacy campaigns. You may argue that this severed ordinary people's connection to Ottoman documents, a real cost. But a clear, shared, easily learned writing standard substantially lowered the threshold for literacy. Standards really can accomplish things.

\textbf{Voice two: ``the standard'' often ratifies the winner after the event.}

The dialect camp's reply is equally complete. First, linguistics, the study of how language works, established long ago that a dialect is not ``incorrect standard language''. Cantonese has a complete, rigorous sound system; Southern Min has a full set of grammatical rules. They are relatives that evolved alongside Putonghua, not defective versions of it. ``Dialects are nonstandard'' has no scientific force in this sense: it is a social verdict dressed as science. Second, the standard's origin cannot withstand questioning. French became ``standard French'' because Paris won. Standard English pronunciation came from southeastern England because the court, Oxford, Cambridge, and the BBC were there. A standard language's history reveals not ``survival of the fittest'' but ``winner takes all'': political victors have their accents retrospectively crowned ``correct''. Third, real people bear the costs: wooden signs, punitive standing, ridicule, rejection at interviews because of accents, and children telling grandparents, ``You say it wrong''. Everyone can count standardisation's dividends; particular groups silently pay its bills.

After hearing both fully, you find they are fighting on different planes. Standard advocates speak of \textbf{function}: a shared language is useful, efficient, and opens opportunities. Dialect advocates speak of \textbf{justice}: why does this useful system take some people's speech as its template while others suffer the pain of changing their voices? One asks ``Is it worthwhile?''; the other, ``By what right?'' The difficulty is that both must be answered together.

\subsection{Honorifics: Putting Status into Grammar}
``Whose speech counts as correct?'' is only the first layer. A finer one governs \textbf{who should speak to whom in what way}.

Japanese has a developed honorific system. The verb for ``eat'' takes completely different forms depending on whether you address a friend, an elder, or a business client. Service industries even have ``manual honorifics'', with every sentence employees address to customers prescribed by company handbooks. Javanese goes further: the language has several levels, with different vocabularies for addressing servants and aristocrats. Opening your mouth first declares both parties' status. Chinese seems to lack such elaborate grammatical machinery, but the distinction between informal ni and respectful nin, the prohibition on addressing elders by personal name, and the ritual of ``After you'' at a meal do the same work.

Consider forms of address too. How people ``should'' be addressed is always political. In the early Chinese Republic, xiansheng and nüshi, ``Mr'' and ``Ms'', displaced laoye and daren, ``Master'' and ``Your Honour''. That announced a new set of relationships, not merely new words. In the later twentieth century, women in the English-speaking world secured ``Ms'', unlike ``Miss'' or ``Mrs'' a title that did not disclose marital status. Why should men use ``Mr'' everywhere while women first announce whether they are married? Behind the invention of a word stood a whole group claiming the right not to have private affairs reported for them.

\subsection{Words You Cannot Say}
Another kind of rule works backwards. It prescribes not what you must say but what you \textbf{must not} say.

Ancient China had naming taboos: nobody could use the emperor's personal name. Emperor Wen of Han was Liu Heng, so the legendary Heng'e became Chang'e. Emperor Taizong of Tang was Li Shimin; official documents had to avoid even min, ``people'', and Tang writers often shortened Guanshiyin to Guanyin. The same word could be writable yesterday and criminal today. The word had not changed; the person sitting above it had. Every society also has softer taboos. Death becomes ``gone'' or ``passed on''; using the toilet becomes ``excusing oneself''. Some things cannot be said before elders, others before children. Taboo vocabulary maps language's minefields, whose boundaries always trace the distribution of power and feeling.

Language rules concern far more than subjects, verbs, and objects. They prescribe standards, ranks, forms of address, and forbidden ground. Together these almost form a social constitution written in words. The difficulty is that its legislators were never elected.

\section[Thinking Bridge: Three Questions]{Thinking Bridge: Turn Correctness into Three Questions}
Next time someone says ``That expression is wrong'' or ``That word is nonstandard'', do not rush to apologise or object. Use a portable tool: split the vague judgement of correctness into three answerable questions.

\textbf{Question one: whose grammar does it violate?}

Every language variety has internal rules. Jinzhao falls outside standard Putonghua usage but fully fits Wu Chinese. In the English of some African American communities, ``He be working'' means ``he keeps working''; this use of be follows rigorous rules rather than random choice. ``Ungrammatical'' therefore requires a completed reference: against \textbf{which} grammar? If the answer is ``the grammar required in examinations'', the issue is not a flaw in the language itself but a mismatch with a particular setting's rules. Making that reference explicit immediately defuses many arguments.

\textbf{Question two: in what setting, and addressed to whom?}

The same sentence changes character with its setting. ``Have you eaten?'' is a greeting in an alley and sarcasm in an interrogation room. Nobody minds shorts when you take out the rubbish; wearing them to receive an award can be a disaster. Language works similarly. Dialect conveys closeness at home and risk in a contract; internet slang builds rapport in a group chat and loses points in a job application. Linguists call this \textbf{register}. People should have several sets of linguistic equipment, just as a wardrobe holds both pyjamas and formal clothes. The question should never be ``Which outfit is correct?'' but ``Does this outfit fit this occasion?'' Someone who can wear only formal clothes and someone who owns only pyjamas share the same limitation.

\textbf{Question three: who benefits from this rule, and who pays?}

This is the sharpest question. Any linguistic rule---standard broadcast pronunciation, an official document format, a forbidden word---can be asked whom it admits and whom it blocks. Examinations accepting only standard written language favour big-city children immersed in it from infancy; children in dialect regions, rural areas, and migrant families pay. Requiring honorifics in formal settings benefits upper groups familiar with them and burdens untrained newcomers. This does not declare all rules conspiracies. Many benefits genuinely extend across society: shared medicine instructions have saved countless lives. It simply requires us to \textbf{calculate the accounts and lay them on the table}, instead of pretending rules are natural laws descending from heaven and benefiting everyone equally.

Together, the questions give you a ``language-rule banknote detector''. Test whether a note marked ``correct'' is currency accepted throughout society or a token printed by a small circle for itself.

\section{Xunzi on Names: An Ancient Answer}
Now read a short primary passage. You might suppose that ``names are conventional'' is a modern discovery. Yet more than two thousand years ago, the Warring States thinker Xunzi wrote ``Rectifying Names'', examining the matter more thoroughly than many people today. One passage says:

\begin{quote}
Names have no inherent appropriateness. They are assigned by agreement; when agreement becomes customary, they are called appropriate. What departs from the agreement is called inappropriate.
\end{quote}
In essence, no name is naturally ``suitable''. People agree to call something by it; once that agreement becomes habitual, it counts as suitable. Violating the agreement counts as unsuitable.

Take this passage apart and weigh it. It has at least four substantial parts.

First, Xunzi dismantles ``innate correctness'' in one sentence. No word is the universe's ordained answer. A table need not be called ``table'' if society originally agreed on something else. When someone calls an expression ``naturally correct'' or ``always so'', invite them to speak with Xunzi.

Second, he blocks the opposite extreme. If names are agreed, may I change them however I please? No. ``What departs from the agreement is called inappropriate.'' Once agreement exists, an individual cannot unilaterally break it. Inventing names only you understand is not avant-garde but a broken connection. Language is public property, so neither an emperor nor each individual can privately monopolise it.

Third, most subtly, ``agreement becoming customary'' hides an unstated difficulty: \textbf{who participated in the agreement?} Everyone at the market contributes to its conventions. But were ordinary people present when court documents, examination standards, and dictionaries established theirs? Xunzi described agreement's power without investigating its procedure. The gap he left, ``Who has the authority to decide?'', remains our central question two thousand years later.

Fourth, his motive in ``rectifying names'' was political. He feared confused names would confuse government orders, rewards, and punishments. Names concern more than linguists. Anyone wishing to govern society eventually discovers that governing language is a shortcut to governing minds.

Calling Xunzi as a witness does not mean he was entirely right. It shows that today's online arguments over what a word means belong to a meeting underway for at least two millennia. Its agenda has never changed: who gets to decide names?

\section[Why Does Standard Speech Pay?]{Why Does Standard Speech Pay? Three Explanations}
Return to the Welsh classroom: English-speaking children advanced; Welsh-speaking children stayed behind. \textbf{Why} did this happen? Separate four things. What happened---opportunities for standard-language speakers, costs for dialect speakers---belongs to evidence. Why it happened invites competing explanations. How people then understood it---the report's description of Welsh as ``backward''---is the participants' account. How we evaluate it is a value judgement. Here we compare the second layer.

\textbf{Explanation one, functional: a shared language lowers the cost of cooperation.}

On this account a standard language wins like a shared railway gauge. Whose gauge it is matters less than having one. Governments need administration, armies commands, markets contracts, and schools textbooks. Every extra language raises transaction costs. English won in Wales not because it was beautiful but because it connected to a larger market. Simple and powerful, this explains why modernisation worldwide accompanies linguistic unification. Its limitation is counting totals without shares. Costs fell, but the \textbf{price} of lowering them landed disproportionately on dialect speakers. Why, moreover, does ``unification'' always mean the capital's language rather than a compromise in which everyone gives ground? Function does not explain that direction.

\textbf{Explanation two, political: standard language is a nation-building project.}

Here unification is not a spontaneous market outcome but a deliberate state construction. Since 1635, the Académie française has been charged with compiling a dictionary and judging usage. Behind modern official languages, national-language campaigns, and writing reforms almost always stands the state. Its reasons are understandable: a population reading newspapers and laws and singing an anthem in one language is easier to identify, mobilise, and govern. This fills the functional account's gap. It explains the standard's \textbf{direction}, why capitals and courts win, and why coercive devices such as wooden signs are widespread. Its limit is that attributing everything to the state cannot explain the frequent absence of compulsion. Welsh parents eagerly sought English lessons; Chinese families eagerly trained children in Putonghua. The state did not force them at gunpoint; they queued themselves.

\textbf{Explanation three, choice: people actively exchange accents for prospects.}

This moves the spotlight from those applying pressure to those choosing. Dialect speakers are not merely passive victims but clear-eyed calculators. They see the pay packets, positions, and dignity connected to the standard and actively change speech, choose Putonghua-speaking kindergartens, and switch to a ``work accent'' on the phone. The account restores people's judgement and initiative and explains why many dialects fade without repression, simply losing speakers. But a question exposes its danger. At a table where every good card goes to standard pronunciation, how voluntary is ``voluntarily changing one's accent''? Describing structural pressure as individual choice gives inequality a comfortable hiding place.

Each explanation holds part of the truth: function reveals benefits, politics reveals direction, and choice reveals human motives. A complete story needs all three. Institutions first raise one language's status, politics; it gains connections to genuine opportunities, function; countless families then make rational but painful changes in speech, choice. The sacrificed language leaves the stage amid applause.

\section[Further Challenge: Language and Power]{Further Challenge: Linguistic Markets and Symbolic Power}
We now need a theory that assembles the pieces. It comes from French sociologist Pierre Bourdieu, whose central idea can be compressed into one sentence: \textbf{linguistic exchange is never merely information exchange; it is also a valuation.}

Bourdieu invites us to imagine a ``linguistic market''. Opening your mouth places a product on the counter: accent, vocabulary, and sentence structure receive an immediate price. Capital-city, elite-school, and broadcast pronunciation are ``hard currency'', redeemable at face value wherever they go. Dialect and foreign accents resemble provincial banknotes: the same meaning is discounted first. Crucially, the speaker does not control the price. Your words may be pearls of wisdom, but the market hears your accent before deciding whether to hear your content.

Abstract as it sounds, the theory has firm empirical foundations. In the 1960s, American linguist William Labov conducted a famous New York experiment. He focused on a sound New Yorkers could pronounce or omit: final r, as in ``fourth floor''. He chose three department stores of distinct status, high-end Saks Fifth Avenue, mid-range Macy's, and lower-end Klein's. Posing as a customer, he asked employees where women's clothing was, prompting ``fourth floor''. The pattern looked ruler-straight: the higher-status the store, the greater the rate of r pronunciation. When he asked staff to repeat themselves, drawing closer attention to their speech, everyone's r use rose sharply, particularly at the middle store, even exceeding the high-end store's level.

The experiment neatly demonstrates the theory's deepest point: \textbf{the linguistic market's most effective regulator lives not in the police station or education ministry but on everyone's own tongue.} Staff knew which pronunciation was ``valuable'' and raised their prices when watched. Middle-store staff raised them most because they were most anxious to climb. No signs or fines were required: people monitored and punished themselves. Bourdieu calls this invisible domination \textbf{symbolic power}. It does not command; it makes certain speech seem ``superior'', and you chase it yourself, looking down on your original speech as you run.

Now consider a counterexample that is easy to misread. You might conclude that because stigma attaches to words, replacing or banning every stigmatised term would solve discrimination. Look at reality. Schools replace ``bad students'' with ``lagging students'', then ``students with learning difficulties''. Each change is kindly intended, but within a few years the new term acquires the old flavour, and children still hear its sting. Linguists call this the ``euphemism treadmill'': words run while discrimination follows a step behind. This does not mean renaming is useless; it is often a group's first step towards dignity. Rather, \textbf{words are where the stain appears, not the stain itself.} Cleaning words without changing attitudes is like wiping a mirror to remove dirt from your face.

This pushes Bourdieu's account further. If linguistic markets automatically reprice new words, what must change is not the merchandise on the counter but the pricing mechanism: who sits in the pricing seat, and by what right?

\section{Today's Language Battlegrounds}
This ancient problem is opening several fronts around you right now.

\textbf{Front one: an accent becomes part of a CV.} During a recruitment call, an accent may change the job the listener imagines for you. In short-video comments, everyone understands the implication after ``You can tell where they're from by the accent''. Meanwhile AI voice assistants, navigation, and customer service all speak flawless standard pronunciation. Machines become correctness's most tireless patrol. As children grow up addressing standard-speaking machines, dialect retreat can only accelerate. This loss needs no wooden sign.

\textbf{Front two: dialects fight back.} On those same platforms, dialect jokes, rap, and dubbed films can win millions of likes. People far from home recognise kin in the comments: ``That accent makes you one of us.'' Language does not merely rank people; it creates belonging. Accents are group passwords. One home-region phrase can make strangers feel like family in a distant restaurant. What the standard-language market discounts is hard currency in the emotional market. Language hurts and excludes, but also gathers people for warmth. These are always two faces of one coin.

\textbf{Front three: words conceal reality.} Examine familiar expressions with this chapter's tools: layoffs become ``optimisation'' or ``graduation'', losses ``negative growth'', reduced benefits ``compensation restructuring''. In ``Politics and the English Language'', Orwell criticised such usage, arguing essentially that vague official abstractions make lies sound sincere and killing respectable. Change the word and the picture disappears. ``Layoff'' contains an actual person carrying a box out of a building; ``optimisation'' contains nobody. Concealment does not require silencing, only a new description without people in it.

\textbf{Front four: naming claims rights.} Naming can also attack. Before the 1970s, English lacked ``sexual harassment''. Office incidents were merely ``jokes'' or ``the boss taking an interest in you''. Feminists coined the term and brought it into law, making that harm something reportable, actionable, and countable. In Chinese, ``domestic violence'' replacing ``family matters'' similarly brought violence behind closed doors into legal light. An unnamed harm resembles a case that cannot be filed. Xunzi described agreement becoming custom; these histories show that \textbf{agreements can be fought for, and the weak can legislate.}

\textbf{Front five: battles over address continue.} Xiaojie, ``Miss'', shifts from respect towards ambiguity; tongzhi, ``comrade'', changes connotations repeatedly; proposals to rename ``migrant workers'' as ``new urban residents'' recur. Online, a label such as ``little pink'', ``fifty steps'', or ``keyboard warrior'' can convict a stranger in three seconds. Labels are violent because they skip the entire process of knowing someone. Once the word is attached, you need not look at the person.

\textbf{Front six: politeness and freedom pull against each other.} This is today's hardest front. One side says using people's preferred terms and avoiding hurtful expressions is basic courtesy, like removing shoes indoors, not oppression. The other says that once forbidden speech can become law or regulation, today's kindness becomes tomorrow's censorship, and freedom shrinks through repeated acts ``for your own good''. Separate facts and judgements first. The fact is that both fears have materialised: groups have gained overdue dignity through changed names, and societies have suppressed legitimate discussion in politeness's name. The judgement is that no rule draws the boundary once and for all. It must be renegotiated in particular conflicts. Uncomfortable, yes, but honest answers often are.

\section[For You: A New Name, a New World?]{For You: Can a New Name Bring a New World?}
Return to the exercise book.

If you could change the rules, what would you do? Option one: allow everything in compositions, including jinzhao and jiaoguan, and remove language-standard marks from examinations. Option two: keep things as they are, circling whatever needs circling; rules are rules. Option three: retain standards but bring dialects into class, teaching both sets of equipment, pyjamas as well as formal clothes.

Before choosing, review the accounts. You have the standard camp's full case: efficiency, fairness, opportunity. Turkish literacy campaigns and medicine instructions are not trivial. You also have the dialect camp's case: standards ratify winners, and signs and accent discrimination are no illusion. You know linguistic markets automatically price people, while seats at the pricing table are never allocated by lot. You know the euphemism treadmill: a new name does not solve the root problem, yet refusing to rename may prevent even a first step. Xunzi has supplied the oldest, sharpest reminder: names have no inherent suitability; everything depends on agreement, and someone is always absent from that table.

Here is the real question: \textbf{if ``correct language'' must be a convention, should we fight for control of its making or broaden participation? Should every variety have a chance to become the standard, or should being ``standard'' itself become less valuable---less able to determine a child's marks, an applicant's prospects, or an older person's dignity at market?}

Can a new name bring a new world? Give your answer and your reasons. Remember that every word you choose while answering already participates in the agreement.
